\chapter{On the Mac, while it is online}
\label{ch:i-mac}

Everything that needs the internet happens here, before the Mac joins the robot's Wi-Fi (it has
no internet). At the end of this chapter the Mac holds the robot's sources in one archive, the
app's APK and a NETRA that runs offline.

\section{Tools}

\begin{lstlisting}[language=bash]
brew install --cask docker android-platform-tools   # Docker Desktop, adb
brew install gh && gh auth login                    # skip what you have already
\end{lstlisting}

Start Docker Desktop once and give it at least 4 GB of memory (Settings › Resources).

\section{The robot's sources in one archive}
\label{sec:i-pack}

\file{tools/pack_sources.sh} clones x30\_ros, every repository of \file{x30.repos} with their
submodules (FAST-LIO with its ikd-Tree, the Livox driver, the HIKMICRO and Hertzino drivers) and
the Livox SDK, and packs them into one archive of about 5 MB, without Git history:

\begin{lstlisting}[language=bash]
mkdir -p ~/x30 && cd ~/x30
git clone git@github.com:DarmawanRobotics/x30_ros.git && cd x30_ros
GIT_CONFIG_COUNT=1 GIT_CONFIG_KEY_0=url.git@github.com:.insteadOf \
  GIT_CONFIG_VALUE_0=https://github.com/ \
  bash tools/pack_sources.sh ~/x30/x30_sources.tar.gz
\end{lstlisting}

The three \code{GIT_CONFIG_*} variables make every clone of this one command go over SSH with
your GitHub key, so private driver repositories come along too; nothing in your Git
configuration changes. \file{src/SOURCES.txt} in the archive lists the commit of every
repository.

\section{The app}

\begin{lstlisting}[language=bash]
cd ~/x30 && gh release download -R DarmawanRobotics/X30Controller -p 'x30-controller-*.apk'
\end{lstlisting}

Every release carries Skydroid's SDK, so it reads the H12's own sticks; on any other Android
device the app drives with touch sticks.

\section{NETRA}
\label{sec:i-netra-build}

NETRA builds its images and downloads its AI model once; afterwards it starts without internet.

\begin{lstlisting}[language=bash]
cd ~/x30 && git clone --recursive git@github.com:DarmawanWeb/Netra.git && cd Netra
scripts/local.sh setup     # checks, .env with generated passwords, Ollama and the model
scripts/local.sh up        # builds and starts everything (a few minutes the first time)
\end{lstlisting}

Open \url{http://localhost} and log in as \code{admin} with the password \cmd{setup} printed
(\code{NETRA_ADMIN_PASSWORD} in \file{.env}).

\subsection{License}

Without a license NETRA runs its core with one robot and drops the inspection uploads of the
robot. The X30 needs the \textbf{substation} module:

\begin{enumerate}
  \item Read the installation ID in \menupath{Settings › License} (e.g.\
    \code{NETRA-7KQ3-M9XD-2WPA}).
  \item Issue the license in the License Manager, or from a checkout of \code{netra-fleet-be}:
\end{enumerate}

\begin{lstlisting}[language=bash]
go run ./cmd/netra-license issue -key ~/netra-keys/ezra.key \
  -client "Gardu X30" -edition substation -tier pro -robots 1 -days 365 \
  -installation NETRA-7KQ3-M9XD-2WPA -out x30.lic
\end{lstlisting}

\begin{enumerate}[resume]
  \item Paste it in \menupath{Settings › License}, or copy it to \file{license/netra.lic} and run
    \cmd{scripts/local.sh up}.
\end{enumerate}

The substation edition draws robots on a GPS basemap. The X30 has no GPS: it appears there
through the site's reference points (\menupath{Site map}). For a drawing of the substation
instead, issue the license with \cmd{-map floorplan}.

\subsection{The robot}

\begin{lstlisting}[language=bash]
scripts/local.sh robot x30-01 "X30 01"
\end{lstlisting}

It prints \code{ROBOT_ID}, \code{CC_URL} and \code{CC_TOKEN}: keep the token, it becomes
\code{X30_NETRA_TOKEN} on the robot (Chapter~\ref{ch:i-netra}). A robot ID is 2 to 32 lower-case
letters, digits, \code{-} and \code{_}.

\section{The handheld's token}

The H12 gateway accepts only a handheld that knows its shared secret. Make one and keep it with
NETRA's:

\begin{lstlisting}[language=bash]
openssl rand -hex 16
\end{lstlisting}

\begin{tipbox}
Before you leave: open \url{http://localhost} once with Wi-Fi off. Everything NETRA needs is
then on the Mac.
\end{tipbox}
